\subsection{Table of Groups of Small Order}

\begin{exercise}[5.3.1]
    Prove that $D_{16}, Z_2 \times D_8, Z_2 \times Q_8, Z_4 * D_8, QD_{16}$, and $M$ are non-isomorphic, non-abelian groups of order 16 (where $Z_4 * D_8$ is described in \exref{5.1.12} and $QD_{16}$ and $M$ are described in \exref{2.5.11} and \exref{2.5.14} respectively).
\end{exercise}

\begin{sol}
    We do a ``decision tree''-like process to show that these are non-isomorphic:
    \begin{itemize}
        \item $D_{16}$
        \begin{itemize}
            \item $Z_2 \times D_8$ has no element of order 8, while $D_{16}$ does.
            \item $Z_2 \times Q_8$ has no element of order 8, while $D_{16}$ does.
            \item $Z_4 * D_8$ has no elements of order 8, since quotienting preserves only divisibility of orders from $Z_4 \times D_8$, which has no elements of order 8, while $D_{16}$ does.
            \item In $QD_{16}$, we have $|\sigma^4| = 2$, and $|\sigma^k\tau| = 2$ when
            \[\sigma^k\tau\sigma^k\tau = \sigma^k\sigma^{3k}\tau\tau = \sigma^{4k},\]
            hence $\sigma^{4k} = 1$ when $k$ is even, hence there are only 5 elements of order 2 in $QD_{16}$, while there are 8 reflections of order 2, plus the central rotation $r^4$, for a total of 9 elements of order 2 in $D_{16}$.
            \item $M$ has $|u| = 2$, with
            \[uv^kuv^k = uuv^{5k}v^k = v^{6k},\]
            so $v^{6k} = 1$ when $8 \mid 6k$, or $4 \mid 3k$, or $4 \mid k$. Then only $k = 4$ works, giving $M$ a total of 3 elements of order 2.
        \end{itemize}
        \item $Z_2 \times D_8$
        \begin{itemize}
            \item $Z_2 \times Q_8$ does not have the same number of elements of order 2 as $Z_2 \times D_8$ does, since $D_8$ has 5 elements of order 2, while $Q_8$ has only 1 element of order 2.
            \item Note that $\bar{(x^k, q)} \in Z_4 * Q_8 \cong Z_4 * D_8$ has order 2 when $(x^k, q)^2 = (1, 1)$ or $(x^2, -1)$. Since we already identify $x^2$ with $-1$ in the quotient, then the case of $q^2 = 1$ and $2i \equiv 0 \bmod 4$ gives the same element, hence only 1 element of order 2. The other case of $q^2 = -1$ and $2i \equiv 2 \bmod 4$, or $i$ odd and $q = Q_8 \setminus Z(Q_8)$ gives 6 elements of order 2, for a total of 7 elements of order 2 in $Z_4 * D_8$, while $Z_2 \times D_8$ has 11 elements of order 2.
            \item $QD_{16}$ contains elements of order 8, while $Z_2 \times D_8$ does not.
            \item $M$ contains elements of order 8, while $Z_2 \times D_8$ does not.
        \end{itemize}
        \item $Z_2 \times Q_8$
        \begin{itemize}
            \item Note that $Z_2 \times Q_8$ only has 3 elements of order 2, while $Z_4 * D_8$ has 7 elements of order 2.
            \item $QD_{16}$ contains elements of order 8, while $Z_2 \times Q_8$ does not.
            \item $M$ contains elements of order 8, while $Z_2 \times Q_8$ does not.
        \end{itemize}
        \item $Z_4 * D_8$
        \begin{itemize}
            \item $QD_{16}$ contains elements of order 8, while $Z_4 * D_8$ does not.
            \item $M$ contains elements of order 8, while $Z_4 * D_8$ does not.
        \end{itemize}
        \item $QD_{16}$
        \begin{itemize}
            \item $M$ has only 3 elements of order 2, while $QD_{16}$ has 5 elements of order 2. \qh
        \end{itemize}
    \end{itemize}
\end{sol}